% !TeX program = xelatex
% 源：教学课件/A往年高一上学期课件/第4章/4.1.1 有理指数幂.pptx
\documentclass[aspectratio=169,10pt]{beamer}
\usepackage[UTF8,fontset=fandol]{ctex}
\IfFontExistsTF{Microsoft YaHei}{\setCJKsansfont{Microsoft YaHei}}{}
\usetheme{GaoZhongYuan}
% 教学内容层；原模板主题、图标及校园素材保持原样。
\usepackage{amsmath,amssymb,booktabs,array,pgfplots}
\pgfplotsset{compat=1.18}
\definecolor{lessonblue}{HTML}{225EA8}
\definecolor{lessonred}{HTML}{BA312D}
\definecolor{lessonpurple}{HTML}{82589F}
\setbeamerfont{normal text}{size=\large}
\setbeamerfont{block title}{size=\normalsize}
\setbeamerfont{block body}{size=\normalsize}
\setbeamertemplate{itemize item}{\color{campusgreen}\small$\bullet$}
\setbeamertemplate{enumerate item}{\color{campusgreen}\insertenumlabel.}
\setlength{\parskip}{.45em}
\setlength{\abovedisplayskip}{.5em}
\setlength{\belowdisplayskip}{.5em}
\newcommand{\R}{\mathbb R}
\newcommand{\Q}{\mathbb Q}
\newcommand{\ans}[1]{{\color{lessonred}#1}}
\newcommand{\key}[1]{{\color{campusgreen}\bfseries #1}}
\newcommand{\reason}[1]{{\color{lessonblue}#1}}
\newcommand{\gap}{\par\vspace{.6em}}
\pgfplotsset{lessonaxis/.style={axis lines=middle,axis line style={->},
 xlabel={$x$},ylabel={$y$},
 tick label style={font=\scriptsize},label style={font=\small},
 width=7.1cm,height=5.1cm,samples=150,clip=true,unbounded coords=jump,
 every axis plot/.append style={thick},
 legend style={font=\scriptsize,draw=none,fill=white,fill opacity=.9,text opacity=1},
 legend cell align=left}}
% 五个基本幂函数，均为可编辑矢量曲线。
\newcommand{\fivepowers}{%
\begin{tikzpicture}\begin{axis}[lessonaxis,xmin=-2.1,xmax=3.2,ymin=-2.1,ymax=3.6,
 xtick={-2,-1,1,2,3},ytick={-2,-1,1,2,3},legend pos=north west]
 \addplot[black,domain=-2:3.1]{x};\addlegendentry{$y=x$}
 \addplot[lessonblue,domain=-1.85:1.85]{x^2};\addlegendentry{$y=x^2$}
 \addplot[lessonred,domain=-1.3:1.5]{x^3};\addlegendentry{$y=x^3$}
 \addplot[campusgreen,domain=0:3.1]{sqrt(x)};\addlegendentry{$y=\sqrt{x}$}
 \addplot[lessonpurple,domain=.285:3.1]{1/x};\addlegendentry{$y=1/x$}
 \addplot[lessonpurple,domain=-2:-.48,forget plot]{1/x};
 \addplot[densely dashed,gray,forget plot] coordinates {(0,1)(1,1)(1,0)};
 \addplot[only marks,mark=*,mark size=1.3pt,forget plot] coordinates {(1,1)};
\end{axis}\end{tikzpicture}}
% 示意值只用于表达 0<q<m<1，未给原题补入数值条件。
\newcommand{\powerorder}{%
\begin{tikzpicture}\begin{axis}[lessonaxis,xmin=0,xmax=3.2,ymin=0,ymax=4,
 xtick={1,2,3},ytick={1,2,3},legend pos=north west]
 \addplot[lessonred,domain=0:1.95]{x^2};\addlegendentry{$y=x^p$}
 \addplot[campusgreen,domain=0:3.1]{x^.6};\addlegendentry{$y=x^m$}
 \addplot[lessonblue,domain=0:3.1]{x^.3};\addlegendentry{$y=x^q$}
 \addplot[lessonpurple,domain=.25:3.1]{1/x};\addlegendentry{$y=x^n$}
 \addplot[densely dashed,gray,forget plot] coordinates {(0,1)(1,1)(1,0)};
\end{axis}\end{tikzpicture}}
\newcommand{\expplot}[1][2]{%
\begin{tikzpicture}\begin{axis}[lessonaxis,xmin=-2.5,xmax=2.5,ymin=0,ymax=5,
 xtick={-2,-1,1,2},ytick={1,2,3,4},legend pos=north west]
 \addplot[lessonred,domain=-2.4:2.3]{2^x};\addlegendentry{$y=2^x$}
 \ifnum#1>1
 \addplot[lessonblue,domain=-2.3:2.4]{2^(-x)};\addlegendentry{$y=(1/2)^x$}
 \fi
 \ifnum#1>2
 \addplot[lessonpurple,domain=-2.4:1.45]{3^x};\addlegendentry{$y=3^x$}
 \addplot[campusgreen,domain=-1.45:2.4]{3^(-x)};\addlegendentry{$y=(1/3)^x$}
 \fi
 \addplot[only marks,mark=*,mark size=1.4pt,black,forget plot] coordinates {(0,1)};
\end{axis}\end{tikzpicture}}
% 四个选项依据原图特征重绘，B 为实立方根。
\newcommand{\rootchoices}{%
\foreach \letter/\kind in {A/1,B/2,C/3,D/4}{%
\begin{minipage}[t]{.24\textwidth}\centering
\begin{tikzpicture}\begin{axis}[lessonaxis,width=3.5cm,height=3.3cm,
 xmin=-1.6,xmax=1.7,ymin=-1.5,ymax=1.8,xtick={1},ytick={1},samples=201]
 \addplot[dashed,gray,domain=-1.4:1.6]{x};
 \ifnum\kind=1 \addplot[black,domain=-1.5:1.5]{tanh(x)};\fi
 \ifnum\kind=2 \addplot[black,domain=-1.5:1.5]{sign(x)*abs(x)^(1/3)};\fi
 \ifnum\kind=3 \addplot[black,domain=-1.12:1.2]{x^3};\fi
 \ifnum\kind=4 \addplot[black,domain=-1.5:1.5]{sign(x+.4)*abs(x+.4)^(1/3)};\fi
\end{axis}\end{tikzpicture}\\[-.5em]\letter
\end{minipage}}\par}

\title{4.1.1 有理指数幂}
\author{深圳中学高中园}
\date{}
\begin{document}
\campusmaketitle

% 原第2页；修正平方根与正边长的区别。
\begin{frame}[t]{任务一：从乘方到开方}
 正方形的边长为 $a$，面积为 $a^2$；立方体的棱长为 $a$，体积为 $a^3$。
 \gap
 \begin{itemize}
  \item 正方形面积为 $4$ 时，边长为 $2$。
  \item 正方形面积为 $5$ 时，边长为 $\sqrt5$。
 \end{itemize}
 \pause
 \key{已知幂与指数，怎样求底数？}
 \[x^2=5\quad\Longrightarrow\quad x=\pm\sqrt5.\]
 作为\key{边长}时取正值 $\sqrt5$。能否把它写成 $5$ 的某次幂？
\end{frame}

% 原第3页
\begin{frame}[t]{任务二：\texorpdfstring{$n$}{n} 次方根}
 \renewcommand{\arraystretch}{1.7}
 \centering
 \begin{tabular}{lll}
  \toprule 方程 & 名称 & 实数解\\\midrule
  $x^2=a\ (a>0)$ & $a$ 的平方根 & $x=\pm\sqrt a$\\
  $x^3=a$ & $a$ 的立方根 & $x=\sqrt[3]a$\\
  $x^4=a\ (a>0)$ & $a$ 的四次方根 & $x=\pm\sqrt[4]a$\\\bottomrule
 \end{tabular}
 \pause
 \begin{block}{定义}
  如果 $x^n=a$，那么 $x$ 叫做 $a$ 的 $n$ 次方根。\\
  其中 $n\in\mathbb N$，且 $n\ge2$。本节在实数范围内讨论。
 \end{block}
\end{frame}

% 原第4页
\begin{frame}[t]{任务二：方根的存在性与个数}
 \renewcommand{\arraystretch}{1.65}
 \centering
 \begin{tabular}{c|c|c}
  \toprule 被开方数 & $n$ 为奇数 & $n$ 为偶数\\\midrule
  $a>0$ & 一个：$\sqrt[n]a>0$ & 两个：$\pm\sqrt[n]a$\\
  $a=0$ & \multicolumn{2}{c}{一个：$0$}\\
  $a<0$ & 一个：$\sqrt[n]a<0$ & 没有实数方根\\\bottomrule
 \end{tabular}
 \pause
 \gap
 \begin{itemize}
  \item 奇次方根对任意实数都存在，并且唯一。
  \item 偶次方根需要 $a\ge0$；当 $a>0$ 时有两个。
 \end{itemize}
\end{frame}

% 原第5页
\begin{frame}[t]{任务三：根式的含义}
 \[\boxed{\sqrt[n]{a}}\qquad n\ge2\]
 \begin{itemize}
  \item $n$ 是\key{根指数}，$a$ 是\key{被开方数}。
  \item 当 $n$ 为奇数时，$\sqrt[n]a$ 表示唯一的实 $n$ 次方根。
  \item 当 $n$ 为偶数时，$\sqrt[n]a$ 表示\key{非负}的 $n$ 次方根。
 \end{itemize}
 \pause
 \[\sqrt9=3,\qquad x^2=9\text{ 的解为 }x=\pm3.\]
 \ans{根式的值与方程的全部实数解需要区分。}
\end{frame}

% 原第6页拆为两页，明确有意义的条件。
\begin{frame}[t]{任务四：先开方，再乘方}
 只要 $\sqrt[n]a$ 在实数范围内有意义，就有
 \[\boxed{\left(\sqrt[n]a\right)^n=a}.\]
 \pause
 \begin{itemize}
  \item $n$ 为奇数：$a\in\R$。
  \item $n$ 为偶数：$a\ge0$。
 \end{itemize}
 \gap
 \[\left(\sqrt[3]{-8}\right)^3=-8,\qquad (\sqrt9)^2=9.\]
 \reason{在有意义的前提下，结果始终等于被开方数 $a$。}
\end{frame}

\begin{frame}[t]{任务四：先乘方，再开方}
 对任意 $a\in\R$，
 \[\boxed{\sqrt[n]{a^n}=\begin{cases}a,&n\text{ 为奇数},\\ |a|,&n\text{ 为偶数}.\end{cases}}\]
 \pause
 \[\sqrt[3]{(-2)^3}=-2,\qquad \sqrt{(-2)^2}=|-2|=2.\]
 \gap
 \ans{偶次根式的值非负，不能直接约去根号与偶数次幂。}
\end{frame}

% 原第7页的引入、定义和运算性质分开。
\begin{frame}[t]{任务五：分数指数的引入}
 \[\sqrt2\cdot\sqrt3=\sqrt6,\qquad
   \sqrt[3]{3^2}\cdot\sqrt[3]4=\sqrt[3]{36}.\]
 根指数不同时，例如 $\sqrt2\cdot\sqrt[3]3$，怎样统一表示？
 \pause
 \begin{block}{正数的分数指数幂}
  对 $a>0$，$m,n$ 为正整数且 $n\ge2$，规定
  \[a^{m/n}=\sqrt[n]{a^m},\qquad
    a^{-m/n}=\frac1{a^{m/n}}=\frac1{\sqrt[n]{a^m}}.\]
 \end{block}
 因而 $\sqrt5=5^{1/2}$，$\sqrt2\cdot\sqrt[3]3=2^{1/2}3^{1/3}$。
\end{frame}

\begin{frame}[t]{任务五：零与底数条件}
 \begin{itemize}
  \item $0$ 的正有理数次幂规定为 $0$。
  \item $0$ 的负有理数次幂没有意义。
  \item 本节不定义 $0^0$。
 \end{itemize}
 \pause
 \gap
 \ans{分数指数幂的统一定义与运算法则在正底数条件下使用。}
 \gap
 负数的奇次方根可以有意义，例如 $\sqrt[3]{-8}=-2$；\\
 这不能作为任意底数都能套用分数指数幂法则的依据。
\end{frame}

\begin{frame}[t]{任务六：有理指数幂的运算}
 对 $a>0$，$b>0$，$r,s\in\Q$，有
 \begin{align*}
  a^r a^s&=a^{r+s},& \frac{a^r}{a^s}&=a^{r-s},\\[.7em]
  (a^r)^s&=a^{rs},& (ab)^r&=a^r b^r.
 \end{align*}
 \pause
 \[\left(\frac ab\right)^r=\frac{a^r}{b^r}.\]
 \gap
 \reason{根式先化成分数指数幂，再按指数运算法则计算。}
\end{frame}

% 原第8页
\begin{frame}[t]{例1：根式求值}
 求下列各式的值：
 \begin{align*}
  &(1)\ \sqrt[3]{(-8)^3}; &&(2)\ \sqrt{(-10)^2};\\[.5em]
  &(3)\ \sqrt[4]{(3-\pi)^4}; &&(4)\ \sqrt{(a-b)^2}\quad(a>b).
 \end{align*}
 \pause
 \begin{align*}
  &(1)\ \ans{-8}; &&(2)\ |-10|=\ans{10};\\[.5em]
  &(3)\ |3-\pi|=\ans{\pi-3}; &&(4)\ |a-b|=\ans{a-b}.
 \end{align*}
 \reason{先看根指数的奇偶性，再判断绝对值内部的符号。}
\end{frame}

% 原第9页，保留两种核心算法。
\begin{frame}[t]{例2：分数指数幂求值}
 \[(1)\ 8^{2/3};\qquad (2)\ \left(\frac{16}{81}\right)^{-3/4}.\]
 \pause
 \begin{align*}
  8^{2/3}&=(2^3)^{2/3}=2^2=\ans4,\\[.5em]
  \left(\frac{16}{81}\right)^{-3/4}
   &=\left[\left(\frac23\right)^4\right]^{-3/4}
     =\left(\frac23\right)^{-3}=\ans{\frac{27}8}.
 \end{align*}
 \pause
 也可以先开方：$8^{2/3}=(\sqrt[3]8)^2=4$。\\
 负指数先取倒数：$\left(\frac{16}{81}\right)^{-3/4}=\left(\frac{81}{16}\right)^{3/4}$。
\end{frame}

% 原第10页
\begin{frame}[t]{例3：根式与指数式互化}
 用分数指数幂表示并化简（$a>0$）：
 \[(1)\ a^2\sqrt[3]{a^2};\qquad (2)\ \sqrt{a\sqrt[3]a}.\]
 \pause
 \begin{align*}
  a^2\sqrt[3]{a^2}&=a^2a^{2/3}=\ans{a^{8/3}},\\[.7em]
  \sqrt{a\sqrt[3]a}&=(a\cdot a^{1/3})^{1/2}
                  =(a^{4/3})^{1/2}=\ans{a^{2/3}}.
 \end{align*}
 \reason{相乘时指数相加，幂的乘方时指数相乘。}
\end{frame}

% 原第11页，独立补全答案。
\begin{frame}[t]{例4：含字母的幂运算}
 计算（式中字母均为正数）：
 \[(1)\ (a^{1/3}b^{-5/6})^6;\]
 \[(2)\ (3m^{3/4}n^{3/2})(-4m^{-1/4}n^{1/3})
          \div(-6m^{1/2}n^{5/6}).\]
 \pause
 \begin{align*}
  (1)&\quad a^{6/3}b^{-30/6}=\ans{a^2b^{-5}=\frac{a^2}{b^5}},\\[.5em]
  (2)&\quad\frac{3(-4)}{-6}\,
  m^{\frac34-\frac14-\frac12}n^{\frac32+\frac13-\frac56}
  =\ans{2n}.
 \end{align*}
\end{frame}

% 原第12页(1)(2)，补全原稿未给出的计算。
\begin{frame}[t]{例5：根式综合化简}
 用分数指数幂表示并化简（字母均为正数）：
 \[(1)\ \frac{\sqrt[3]{a^2}\sqrt b}{\sqrt a\sqrt[3]b};\qquad
   (2)\ (\sqrt[3]3-\sqrt{27})\div\sqrt[6]3.\]
 \pause
 \begin{align*}
  (1)&\quad a^{\frac23-\frac12}b^{\frac12-\frac13}
       =\ans{a^{1/6}b^{1/6}},\\[.7em]
  (2)&\quad 3^{\frac13-\frac16}-3^{\frac32-\frac16}
       =\ans{3^{1/6}-3^{4/3}}.
 \end{align*}
 \ans{商对差式运算时，差式中的每一项都要除以除数。}
\end{frame}

% 原第12页(3)，根号范围与原图核对。
\begin{frame}[t]{例5：嵌套根式}
 化简（$x>0$，$y>0$）：
 \[\sqrt[3]{xy^2(\sqrt{xy})^3}.\]
 \pause
 \begin{align*}
  \sqrt[3]{xy^2(\sqrt{xy})^3}
  &=\left[xy^2(xy)^{3/2}\right]^{1/3}\\[.5em]
  &=\left(x^{5/2}y^{7/2}\right)^{1/3}
    =\ans{x^{5/6}y^{7/6}}.
 \end{align*}
 \reason{先明确每层根号的范围，再从内到外改写指数。}
\end{frame}

\begin{frame}[t]{课堂小结}
 \begin{enumerate}
  \item $n$ 次方根：先判断存在性，再判断个数。
  \item 两个基本公式：
   \[(\sqrt[n]a)^n=a,\qquad
     \sqrt[n]{a^n}=\begin{cases}a,&n\text{ 奇},\\|a|,&n\text{ 偶}.\end{cases}\]
  \item 正底数的根式可转为分数指数幂，统一使用指数法则。
 \end{enumerate}
 \gap
 \key{自检：条件写全了吗？偶次根式非负吗？负指数处理正确吗？}
\end{frame}
\campusclosing[谢谢]
\end{document}
